\documentclass[conference]{IEEEtran}
\IEEEoverridecommandlockouts
\usepackage{cite}
\usepackage{amsmath,amssymb,amsfonts}
\usepackage{amsthm}
\usepackage{algorithmic}
\newtheorem{definition}{Definition}
\newtheorem{proposition}{Proposition}
\usepackage{graphicx}
\usepackage{textcomp}
\usepackage{xcolor}
\usepackage[T1]{fontenc}
\usepackage{booktabs}
\usepackage{orcidlink}
\usepackage{multirow}
\usepackage{url}
\usepackage{algorithm}
\usepackage{balance}

% IEEE style: Roman numerals for tables
\renewcommand{\thetable}{\Roman{table}}

\title{Quantum-Inspired Uncertainty Aggregation for Materials Active Learning:\\Parity via Model-Disagreement Coupling, Not Covariance}

\author{
\IEEEauthorblockN{Arnav Kapoor\orcidlink{0009-0007-9818-7908}}
\IEEEauthorblockA{
Indian Institute of Science Education and Research Bhopal, Bhopal 462066, India\\
arnavkapoor23@iiserb.ac.in
}
}

\begin{document}
\maketitle

\begin{abstract}
Acquisition functions in active learning for materials science typically collapse multi-property uncertainty to a scalar score, discarding cross-property correlations that could in principle carry information about sample informativeness under a limited query budget. We formalize a quantum-inspired framework in which each candidate material is encoded as a normalized state in a feature Hilbert space, uncertainty is evaluated through a set of non-commuting Hermitian observables corresponding to distinct physicochemical aspects, and individual variances together with cross-observable covariances are aggregated via complex-valued coupling coefficients into a single, covariance-aware acquisition score, with a proof that it reduces analytically, in the commuting real-coefficient limit, to a standard weighted-variance acquisition rule. We implement this formalism exactly as specified, unit-test it against that classical-limit reduction, and evaluate it against nine standard active-learning baselines on five real regression tasks drawn from the Materials Project ($N\approx 1000$ materials each), with five independent trials, paired significance testing, and an ablation study. \textbf{As originally specified, the formalism underperforms}: it loses on four of five tasks, and an ablation shows the covariance-coupling mechanism it is built around changes performance by an amount indistinguishable from noise ($+0.07\%$ $R^2$ against $\sigma\approx5$--$8\%$). We diagnose the cause: the acquisition score is a function only of a candidate's static position in feature space and never observes the downstream predictor's actual residuals or disagreement, unlike every baseline that wins. Two narrow fixes (domain-informed observable grouping; predictor-importance-weighted encoding) fail to address this. A third, more direct fix -- encoding candidate states from the per-tree predictions of the downstream random forest, i.e.\ genuinely injecting ensemble disagreement into the formalism -- closes the gap: the resulting variant is statistically indistinguishable from the best classical baseline on \emph{all five} real tasks ($p>0.15$ throughout, no significant loss anywhere). However, a further ablation of this fixed variant shows the improvement is attributable entirely to exposure to ensemble disagreement, not to the quantum-inspired machinery: dropping the covariance term or collapsing to a single observable performs as well or better than the full multi-observable, complex-coupled construction on four of five tasks, and the full covariance-aware variant is markedly \emph{unstable} on the highest-variance target (dielectric constant), underperforming even its own single-observable ablation there by a wide margin. We conclude that reaching competitive performance required abandoning the paper's central novelty -- covariance-aware aggregation over non-commuting observables -- in favor of directly measuring model disagreement, which is what the winning classical baselines (QBC, RF-Uncertainty) already do, more simply and more robustly. Separately, we build and verify an actual quantum circuit realization of the formalism (amplitude encoding plus Pauli-decomposed measurement) and confirm it reproduces the classical closed-form exactly; we then characterize its NISQ resource cost -- 528 Pauli measurement terms per variance/covariance quantity, combinatorial growth with the number of observables, and shot/noise sensitivity of the resulting acquisition ranking -- and show that standard qubit-wise-commuting measurement grouping provides a genuine 4.3$\times$ reduction in that cost (12.3$\times$ with full commuting grouping, at added circuit-depth cost). This is an honest feasibility study, not a performance claim.
\end{abstract}

\begin{IEEEkeywords}
active learning, materials discovery, quantum-inspired algorithms, uncertainty quantification, negative result
\end{IEEEkeywords}

\section{Introduction}

The computational and experimental cost of labeling new materials candidates remains a principal bottleneck in accelerated discovery pipelines \cite{wang2023scientific,ren2018accelerated,pyzerknapp2022accelerating}. Active learning (AL) alleviates this burden by sequentially selecting the most informative candidates for evaluation under a fixed labeling budget \cite{settles2009active,cohn1996active}. The central algorithmic ingredient is an acquisition function that scores unlabeled candidates by their expected contribution to model improvement.

Most established acquisition functions -- uncertainty sampling, Query-by-Committee (QBC), Expected Improvement (EI), and diversity-based heuristics -- operate on a single scalar uncertainty estimate \cite{settles2012active,sener2018active,ash2020deep}. In multi-property materials discovery, target quantities such as band gap, formation energy, and elastic response are not independent; they share compositional and structural determinants. This motivates a natural question: can an acquisition function that \emph{explicitly} represents cross-property/cross-observable uncertainty structure -- rather than collapsing it to a scalar before aggregation -- select more informative batches?

We formalize one specific, mathematically principled answer to this question, borrowing the operator language of quantum mechanics as a structured way to combine multiple uncertainty ``observables'' with complex coupling coefficients (Section~\ref{sec:method}). Unlike an earlier version of this manuscript, every number reported below comes from an actual run of an actual implementation of that formalism (\texttt{quantum\_al/operator.py}, unit-tested against the classical-limit proof in Proposition~\ref{prop:classical}), evaluated on genuine Materials Project data, with a fixed, pre-registered protocol, real baselines, and real statistical tests.

\textbf{As originally specified, the answer is no.} The formalism does not beat any of nine standard baselines on four of five real regression tasks. We diagnose why (Section~\ref{sec:improvement}) and, guided by that diagnosis, construct a variant that reaches statistical parity with the best baseline on every task -- but only by directly injecting ensemble disagreement into the state encoding, at which point the formalism's own covariance/multi-observable machinery turns out to contribute nothing beyond that signal, and is actively destabilizing on the noisiest task. We report the full arc, not just the flattering parts.

\subsection*{Contributions}
\begin{enumerate}
    \item A complete, correctly-implemented and unit-tested realization of a covariance-aware, quantum-inspired uncertainty formalism (state encoding, non-commuting Hermitian observables, symmetrized covariance, complex-coefficient acquisition score, and a verified classical-limit reduction).
    \item A rigorous empirical evaluation against nine active-learning baselines on five real Materials Project regression tasks, with paired significance testing, an ablation study isolating each claimed mechanism, and measured runtime/memory. As specified, the formalism underperforms, and none of its claimed mechanisms (covariance coupling, non-commutativity, complex coefficients) produce an effect distinguishable from noise.
    \item A concrete diagnosis (the acquisition score never observes the predictor's residuals) and two narrow fixes that fail to address it (domain-informed observable grouping; predictor-importance-weighted encoding).
    \item A third fix -- coupling the state encoding to the downstream random forest's per-tree predictions -- that reaches parity with the best classical baseline on all five real tasks, together with a further ablation showing that parity is attributable entirely to the ensemble-disagreement coupling, not to the quantum-inspired covariance apparatus, which remains inert or actively harmful even in this improved variant.
    \item An actual quantum circuit realization of the formalism (Section~\ref{sec:hardware}), verified to reproduce the classical closed-form exactly, together with an honest NISQ resource characterization: Pauli measurement overhead, shot-count convergence of the acquisition ranking, and sensitivity to realistic gate noise.
\end{enumerate}

\section{Theoretical Framework}\label{sec:method}

\subsection{State encoding}

Let $x_i \in \mathbb{R}^d$ denote the standardized feature vector of candidate material $M_i$. We embed $x_i$ into a $d$-dimensional complex Hilbert space $\mathcal{H}_d$ via amplitude encoding:

\begin{definition}[Quantum state of a material]\label{def:state}
The quantum state of candidate $M_i$ is
\begin{equation}
|\psi_i\rangle = \sum_{j=1}^{d} \alpha_{ij} |f_j\rangle, \label{eq:state}
\end{equation}
where $\{|f_j\rangle\}_{j=1}^{d}$ is an orthonormal basis corresponding to feature dimensions and $\alpha_{ij} = \sqrt{p_{ij}}\, e^{i\phi_{ij}}$ satisfies $\sum_j |\alpha_{ij}|^2 = 1$.
\end{definition}

We operationalize $p_{ij}$ and $\phi_{ij}$ concretely (the original manuscript left this underspecified): $p_{ij} = \mathrm{softmax}_j(x_{ij}^2)$, so the normalization constraint holds by construction; and $\phi_{ij} = \phi_j$ (constant across candidates for a fixed labeled pool) is $\pi$ times the mean absolute Pearson correlation of feature $j$ with the other $d-1$ features, estimated over the \emph{current labeled set} and recomputed at every AL iteration. This choice makes the phase term of Eq.~\eqref{eq:state} a genuine, if simple, encoding of the pool's correlation structure, and is the interpretation we implement and test throughout.

\subsection{Non-commuting observables}

We define $K$ Hermitian operators $\hat{O}_k : \mathcal{H}_d \to \mathcal{H}_d$,
\begin{equation}
\hat{O}_k = \sum_{i,j} w^{(k)}_{ij} |f_i\rangle\langle f_j|, \qquad \hat{O}_k = \hat{O}_k^\dagger, \label{eq:observable}
\end{equation}
with $K=3$ observables -- structural, electronic, thermodynamic -- each concentrated (via the weight matrix $W^{(k)}$) on a distinct, overlapping subset of feature indices, normalized to unit spectral norm. Non-commutativity ($[\hat{O}_k,\hat{O}_l]\neq 0$) is guaranteed by construction (verified numerically for every run; see Section~\ref{sec:results}).

\subsection{Uncertainty and covariance}

For each observable, the quantum variance of state $|\psi_i\rangle$ is
\begin{equation}
\mathrm{Var}(\hat{O}_k;\psi_i) = \langle\psi_i|\hat{O}_k^2|\psi_i\rangle - \langle\psi_i|\hat{O}_k|\psi_i\rangle^2. \label{eq:variance}
\end{equation}
Cross-observable coupling is captured by the symmetrized covariance
\begin{equation}
\mathrm{Cov}(\hat{O}_k,\hat{O}_l;\psi_i) = \left\langle\psi_i\left|\tfrac{\hat{O}_k\hat{O}_l+\hat{O}_l\hat{O}_k}{2}\right|\psi_i\right\rangle - \langle\psi_i|\hat{O}_k|\psi_i\rangle\langle\psi_i|\hat{O}_l|\psi_i\rangle. \label{eq:covariance}
\end{equation}
Symmetrization guarantees $\mathrm{Cov}\in\mathbb{R}$.

\begin{definition}[Total uncertainty score]\label{def:total}
Given complex coupling coefficients $\{\alpha_k\}_{k=1}^K \subset \mathbb{C}$, the covariance-aware acquisition score for candidate $M_i$ is
\begin{equation}
U_{\mathrm{total}}(M_i) = \sqrt{\sum_{k=1}^K |\alpha_k|^2 \mathrm{Var}(\hat{O}_k;\psi_i) + \sum_{\substack{k,l=1\\k\neq l}}^{K} \mathrm{Re}(\alpha_k^*\alpha_l)\,\mathrm{Cov}(\hat{O}_k,\hat{O}_l;\psi_i)}. \label{eq:utotal}
\end{equation}
\end{definition}

Candidates are ranked by $U_{\mathrm{total}}$ and the top $b$ form the query batch.

\begin{proposition}[Classical reduction]\label{prop:classical}
If (i) all observables commute, $[\hat{O}_k,\hat{O}_l]=0$; (ii) all coefficients are real, $\alpha_k \in \mathbb{R}$; and (iii) the covariance terms are dropped, then
\begin{equation}
U_{\mathrm{total}}^2 \to \sum_{k=1}^K \alpha_k^2\, \mathrm{Var}(\hat{O}_k;\psi_i),
\end{equation}
a standard weighted-variance aggregator over independent observables.
\end{proposition}

\texttt{quantum\_al/operator.py} implements Eqs.~\eqref{eq:state}--\eqref{eq:utotal} directly and includes a numerical self-test verifying (a) non-commutativity of the default $K{=}3$ observable bank and (b) Proposition~\ref{prop:classical} to floating-point precision ($|U_{\mathrm{total}} - \sqrt{\sum_k \alpha_k^2\mathrm{Var}_k}| < 10^{-9}$).

\subsection{Computational complexity}

Every expectation is a quadratic form, $\langle\psi|\hat{O}|\psi\rangle = \alpha^\dagger W \alpha$, costing $O(d^2)$ per observable for dense $W$. For $N$ candidates and $K$ observables, total complexity is $O(NKd^2)$; in our experiments ($N\lesssim 1000$, $K=3$, $d=21$) this is negligible in absolute terms (Section~\ref{sec:runtime}).

\section{Experimental Setup}\label{sec:setup}

\subsection{Datasets}

Unlike the synthetic or partially-synthetic datasets used during earlier development of this codebase, every regression target reported here is a genuine DFT-derived quantity pulled directly from the Materials Project via its public API \cite{jain2013materials}. Table~\ref{tab:dataset} summarizes the five tasks. \emph{Thermal conductivity, claimed in an earlier version of this manuscript, is not included}: the Materials Project does not carry broadly computed thermal-conductivity data, so rather than substitute a synthetic proxy we simply evaluate on the five properties that are genuinely available at scale.

\begin{table}[!t]
\centering
\caption{Dataset summary. All values are real DFT-derived quantities from the Materials Project (queried live; $N=1000$ per task before any downstream filtering).}
\label{tab:dataset}
\begin{tabular}{lccc}
\toprule
Property & $N$ & Range & Source field \\
\midrule
Band gap (eV) & 1000 & 0.00--7.41 & \texttt{band\_gap} \\
Formation energy (eV/atom) & 1000 & $-4.41$--$2.30$ & \texttt{formation\_energy\_per\_atom} \\
Bulk modulus (GPa) & 1000 & 1.8--326.0 & \texttt{bulk\_modulus.vrh} \\
Magnetic moment ($\mu_B$) & 1000 & 0.01--96.0 & \texttt{total\_magnetization} \\
Dielectric constant & 1000 & 1.3--1455.2 & \texttt{e\_total} \\
\bottomrule
\end{tabular}
\end{table}

\textbf{Features} are 21-dimensional (not 100-dimensional, as claimed in an earlier version of this manuscript): six structural/summary descriptors (number of elements, density, volume per atom, site count, energy above hull, space-group number) directly from the Materials Project summary endpoint, plus fifteen composition-derived statistics (mean/std/range, weighted over the constituent elements, of electronegativity, atomic radius, atomic mass, periodic-table row, and periodic-table group). All continuous features are standardized (zero mean, unit variance, fit on the training pool) before state encoding.

\subsection{Baselines}

We compare against nine Tier-1 acquisition strategies, each a real, working implementation (not a placeholder): Uncertainty Sampling (Gaussian-process posterior std), Query-by-Committee (5-model ensemble disagreement), Expected Improvement, Maximum Entropy (random-forest tree-prediction entropy proxy), Diversity Sampling ($k$-means), BADGE, CoreSet (greedy $k$-center), Random-Forest variance, and Random Sampling. We do not claim comparison against Tier-2/3 methods (full multi-output GPs, covariance-calibrated deep ensembles, multi-task predictors) from an earlier version of this manuscript -- those were never actually implemented in this codebase and are left as future work.

\subsection{Protocol}

Five independent trials (seeds 0--4), 70\%/30\% pool-test split, $n_0=50$ initial labeled samples drawn uniformly at random, $T=8$ AL iterations, batch size $b=15$. The downstream predictor is a random forest with 100 trees, retrained from scratch each iteration. Primary metrics are $R^2$ and MAE on the held-out test split.

\section{Results}\label{sec:results}

\subsection{Primary comparison}

Table~\ref{tab:primary} reports final-iteration $R^2$ across all five tasks. The proposed method is \emph{last} of ten on band gap, 6th of 10 on formation energy, 8th of 10 on bulk modulus, and \emph{last} on magnetic moment -- where it is the only method with negative $R^2$, i.e.\ worse than predicting the training mean. It is nominally best on dielectric constant, but every method achieves $R^2\lesssim 0.03$ there; this reflects an intrinsically hard prediction target from the available features, not a genuine advantage.

\begin{table}[!t]
\centering
\caption{Final test $R^2$ (mean $\pm$ std, 5 trials, 8 iterations) on band gap, the task used for statistical testing and ablations. Full five-task table in \texttt{results/primary\_benchmark\_summary.json}.}
\label{tab:primary}
\begin{tabular}{lc}
\toprule
Method & Band gap $R^2$ \\
\midrule
Uncertainty Sampling & $0.5966 \pm 0.0432$ \\
Query by Committee & $0.5958 \pm 0.0419$ \\
Maximum Entropy / RF Uncertainty & $0.5827 \pm 0.0459$ \\
BADGE & $0.5814 \pm 0.0627$ \\
Random Sampling & $0.5806 \pm 0.0667$ \\
CoreSet & $0.5558 \pm 0.0272$ \\
Diversity Sampling & $0.5531 \pm 0.0599$ \\
Expected Improvement & $0.5508 \pm 0.1046$ \\
\textbf{Quantum-Inspired (ours)} & $\mathbf{0.5131 \pm 0.0458}$ \\
\bottomrule
\end{tabular}
\end{table}

\subsection{Statistical testing}

Table~\ref{tab:stats} reports paired two-sided $t$-tests on band gap (quantum vs.\ each baseline, 5 trials). \emph{Every} comparison favors the baseline (mean difference $<0$). Several raw $p$-values fall below 0.05, but after Holm--Bonferroni correction across the nine comparisons, none remain significant at $\alpha=0.05$ -- so we cannot even claim a statistically confident \emph{loss}, only a consistent, non-significant disadvantage. Shapiro--Wilk tests confirm the paired differences are approximately normal in every case ($p>0.05$), validating the paired $t$-test.

\begin{table}[!t]
\centering
\caption{Paired $t$-test results, quantum vs.\ baseline, on band gap (n=5 trials). Positive mean-diff would favor the quantum method; none are positive.}
\label{tab:stats}
\begin{tabular}{lccc}
\toprule
Comparison (vs.\ quantum) & $t$ & $p$ (raw) & $p$ (Holm--Bonf.) \\
\midrule
Uncertainty Sampling & $-4.53$ & 0.011 & 0.073 \\
Query by Committee & $-4.55$ & 0.010 & 0.073 \\
Maximum Entropy & $-4.95$ & 0.008 & 0.070 \\
Expected Improvement & $-1.01$ & 0.371 & 0.410 \\
Diversity Sampling & $-1.51$ & 0.205 & 0.410 \\
\bottomrule
\end{tabular}
\end{table}

\subsection{Ablation study}

Table~\ref{tab:ablation} isolates each mechanism the formalism is built around. \textbf{None of them matter}: removing the covariance term changes $R^2$ by $+0.0007$ (i.e.\ nothing), forcing commuting observables changes it by $+0.0062$, restricting to real coefficients changes it by $-0.0081$, and collapsing to a single observable changes it by $-0.0107$ -- all well within one trial-to-trial standard deviation ($\sigma\approx 0.05$--$0.08$). This directly contradicts the mechanistic narrative of an earlier version of this manuscript, which claimed covariance coupling was ``the primary driver of improvement'' with a $-2.8\%$ $R^2$ drop from its removal; on real data, the measured effect is two orders of magnitude smaller and not distinguishable from noise.

\begin{table}[!t]
\centering
\caption{Ablation on band gap (mean $R^2$, 5 trials; full model $=0.5131$).}
\label{tab:ablation}
\begin{tabular}{lcc}
\toprule
Variant & $R^2$ & $\Delta R^2$ \\
\midrule
Full model & $0.5131 \pm 0.0458$ & --- \\
No covariance ($\mathrm{Cov}=0$) & $0.5137 \pm 0.0763$ & $+0.0007$ \\
Commuting observables only & $0.5192 \pm 0.0588$ & $+0.0062$ \\
Real-only coefficients ($\alpha_k\in\mathbb{R}$) & $0.5050 \pm 0.0742$ & $-0.0081$ \\
Single observable ($K=1$) & $0.5023 \pm 0.0652$ & $-0.0107$ \\
\bottomrule
\end{tabular}
\end{table}

\subsection{Observable sensitivity}\label{sec:obssens}

Sweeping the number of observables on band gap (5 trials each): $K=2$ gives $R^2=0.4759\pm0.0682$, $K=3$ (default) gives $0.5131\pm0.0458$, and $K=6$ gives $0.5122\pm0.0689$. $K=3$ is marginally best among those tested -- directionally consistent with the original claim -- but the method trails every baseline regardless of $K$.

\subsection{Runtime and memory}\label{sec:runtime}

Table~\ref{tab:runtime} reports measured (not simulated) wall-clock time and peak memory on the real band-gap candidate pool ($N=650$, single CPU thread, 3 repeats). The proposed method is fast and lightweight in absolute terms -- faster than QBC and Maximum Entropy -- but this is not a meaningful advantage given it does not win on predictive performance.

\begin{table}[!t]
\centering
\caption{Measured acquisition time and memory at real $N=650$ (mean $\pm$ std, 3 runs, single CPU thread).}
\label{tab:runtime}
\begin{tabular}{lcc}
\toprule
Method & Time (s) & Memory (MB) \\
\midrule
Uncertainty Sampling & $0.081\pm0.004$ & $0.79\pm0.01$ \\
Expected Improvement & $0.083\pm0.001$ & $1.28\pm0.00$ \\
Quantum-Inspired (ours) & $0.369\pm0.002$ & $0.67\pm0.00$ \\
Maximum Entropy & $0.880\pm0.040$ & $1.17\pm0.00$ \\
Query by Committee & $1.539\pm0.019$ & $0.96\pm0.07$ \\
\bottomrule
\end{tabular}
\end{table}

\subsection{Improvement attempt}\label{sec:improvement}

Before concluding, we tested two principled, pre-specified modifications, isolated from each other, to check whether the negative result was an artifact of implementation choices rather than the formalism itself:

\begin{enumerate}
    \item \emph{Domain-informed observable grouping.} The default construction assigns feature indices to observables via arbitrary contiguous thirds, which mixes semantically unrelated columns. We reassigned features to structural/electronic/thermodynamic groups by actual physical meaning (e.g.\ electronegativity and periodic-table group to ``electronic''; atomic radius and density to ``structural''). Result: $R^2=0.4915\pm0.0903$ -- \emph{worse} than the arbitrary grouping.
    \item \emph{Predictor-importance-weighted encoding.} We rescaled each standardized feature by the square root of its importance in a random forest fit on the current labeled set (the same information every winning baseline already exploits) before state encoding, on top of the domain-informed grouping. Result: $R^2=0.5124\pm0.0773$ -- statistically indistinguishable from the unmodified original.
\end{enumerate}

Neither change closes the gap to the best baseline ($R^2=0.5966$). This is evidence that the shortfall is structural rather than a matter of feature-grouping or reweighting details: $U_{\mathrm{total}}(M_i)$ is a function only of candidate $M_i$'s position in a standardized feature space and the labeled pool's correlation structure. It never observes the trained predictor's residuals or disagreement.

\subsection{A third fix: coupling to ensemble disagreement}\label{sec:v3}

The diagnosis above is specific enough to suggest a direct fix, rather than a workaround: make $|\psi_i\rangle$ a function of the downstream predictor's disagreement rather than of the raw material features. Concretely, we fit the same random forest used for evaluation on the current labeled set, collect each candidate's vector of \emph{per-tree predictions} (one scalar per tree), standardize this vector across the candidate pool, and encode \emph{that} as the quantum state via Eq.~\eqref{eq:state} -- trees replace features as the basis, split into 3 arbitrary equal-size groups (``ensemble observables''; there is no structural/electronic/thermodynamic meaning once the basis is tree predictions rather than material descriptors), with the same variance/covariance/complex-coupling machinery of Eqs.~\eqref{eq:variance}--\eqref{eq:utotal} applied unmodified on top.

Table~\ref{tab:v3} reports this variant (\texttt{Quantum-TreeEnsemble}) against each task's best classical baseline, 5 trials, same protocol. \textbf{It reaches statistical parity on all five real tasks} -- no paired comparison approaches significance ($p>0.15$ throughout) -- a qualitative change from the original formalism, which lost significantly on band gap.

\begin{table}[!t]
\centering
\caption{Residual-coupled variant (\texttt{Quantum-TreeEnsemble}) vs.\ each task's best classical baseline, mean $R^2\pm$std, 5 trials, paired $t$-test.}
\label{tab:v3}
\begin{tabular}{lccc}
\toprule
Task & Quantum-TreeEns. & Best baseline & $p$ \\
\midrule
Band gap & $0.590\pm0.066$ & US: $0.597$ & 0.638 \\
Formation energy & $0.821\pm0.027$ & MaxEnt: $0.853$ & 0.161 \\
Bulk modulus & $0.803\pm0.023$ & Divers.: $0.826$ & 0.163 \\
Magnetic moment & $0.323\pm0.068$ & BADGE: $0.418$ & 0.150 \\
Dielectric const. & $-0.450\pm0.724$ & BADGE: $-0.020$ & 0.322 \\
\bottomrule
\end{tabular}
\end{table}

This looks, at first glance, like the positive result the original manuscript claimed. \textbf{It is not evidence for the covariance-aware quantum formalism specifically.} We ran the same ablation as Table~\ref{tab:ablation} on this new variant, across all five tasks (Table~\ref{tab:v3ablation}): dropping the covariance term, or collapsing to a single observable (i.e.\ a bare quadratic form over per-tree predictions, structurally close to the RF-Uncertainty baseline's raw tree-prediction variance), performs \emph{as well or better} than the full multi-observable, complex-coupled construction on four of five tasks. On the fifth and noisiest task (dielectric constant, target range 1.3--1455), the full covariance-aware variant is not just worse but unstable ($R^2=-0.450\pm0.724$, individual trials as low as $-1.86$), while its own single-observable ablation is close to the best baseline ($R^2=-0.160$). Coupling multiple non-commuting observables and complex coefficients on top of a genuinely informative signal did not preserve that signal's value; on the hardest task it actively degraded it.

\begin{table}[!t]
\centering
\caption{Ablation of the residual-coupled variant, mean final $R^2$, 5 trials, all 5 tasks. Bold marks the best of the three per task.}
\label{tab:v3ablation}
\begin{tabular}{lccc}
\toprule
Task & Full (Cov+$K$=3) & No covariance & Single obs.\ ($K$=1) \\
\midrule
Band gap & 0.590 & 0.572 & \textbf{0.614} \\
Formation energy & 0.821 & 0.809 & \textbf{0.844} \\
Bulk modulus & 0.803 & 0.796 & \textbf{0.808} \\
Magnetic moment & \textbf{0.323} & 0.277 & 0.298 \\
Dielectric const. & $-0.450$ & \textbf{$-0.022$} & $-0.160$ \\
\bottomrule
\end{tabular}
\end{table}

The pattern is consistent: whatever value this variant has comes from exposing the acquisition score to ensemble disagreement at all, not from the specific way the formalism aggregates it. We did not search further after this; three independent, pre-specified architectural interventions (Section~\ref{sec:improvement}, and this one) is already enough attempts that continuing to search for a configuration that wins would risk fitting noise in a five-task, five-trial evaluation rather than finding a genuine effect.

\section{Secondary Experiments}

Three further experiments from an earlier version of this manuscript were rerun on the present codebase; we report their status honestly rather than omit them.

\textbf{Discrete classification.} On a synthetic 6-class crystal-system task (not yet wired to the real \texttt{data/crystal\_system.json} labels collected alongside the regression data), final accuracy after 8 iterations was: quantum-margin sampling 89.3\%, entropy sampling 89.7\%, random sampling 79.3\%. Quantum-margin is essentially tied with entropy sampling (entropy marginally ahead) and both clearly beat random. This is a much more modest finding than an earlier claim of ``91\% vs.\ 85\%,'' and does not show the quantum method beating entropy sampling as previously claimed.

\textbf{Transfer learning.} The rerun output for the oxide$\to$sulfide/nitride transfer experiment showed the "transfer" and "train-from-scratch" conditions producing byte-identical $R^2$ trajectories. This indicates an implementation bug (the two code paths are not actually diverging), not a genuine finding, and is flagged here rather than reported as a result. A corrected version is left for future work.

\textbf{Multi-property optimization.} The rerun multi-property active-learning loop produced $R^2$ values near zero or slightly negative across all four properties and all iterations, indicating the loop is not learning effectively in its current synthetic setup. This is inconclusive and flagged for follow-up rather than reported as a positive or negative finding.

\textbf{Spurious-correlation analysis.} This experiment is a deliberately synthetic noise-injection study of a failure mode (not intended to model real materials) and ran successfully; its output is available in \texttt{results/spurious\_correlation.json} for reference but is not load-bearing for the paper's central claim.

\section{Quantum Hardware Realization}\label{sec:hardware}

Everything reported above is quantum-\emph{inspired}: $|\psi_i\rangle$ is a NumPy complex vector and $\hat{O}_k$ a NumPy matrix, manipulated by ordinary linear algebra -- no circuit, qubit, or measurement appears anywhere in Sections~\ref{sec:results}--\ref{sec:v3}. This section builds and characterizes what an actual quantum circuit realization of Eqs.~\eqref{eq:state}--\eqref{eq:utotal} looks like, which the manuscript's original future-work discussion gestured at but never implemented: amplitude encoding for state preparation, and Pauli-decomposed measurement circuits for $\mathrm{Var}$ and $\mathrm{Cov}$. This is not an attempt to show a performance advantage -- it is a resource and feasibility characterization, reported honestly regardless of outcome, implemented in \texttt{quantum\_al/circuit.py} using Qiskit and Qiskit Aer.

\subsection{Correctness: the circuit matches the formalism exactly}

For each candidate, $x_i$ is amplitude-encoded onto $n=\lceil\log_2 d\rceil$ qubits ($n=5$ for $d=21$, zero-padded to dimension 32) via a state-preparation circuit, and each observable $\hat{O}_k$ (correspondingly zero-padded) is decomposed into a sum of Pauli strings, $\hat{O}_k = \sum_p c_p P_p$, via exact trace-based projection. In the noiseless, infinite-precision limit (exact statevector simulation), the resulting circuit-measured $\mathrm{Var}(\hat{O}_k;\psi_i)$ and $\mathrm{Cov}(\hat{O}_k,\hat{O}_l;\psi_i)$ match the classical closed-form values in \texttt{quantum\_al/operator.py} to floating-point precision ($<10^{-9}$; see \texttt{quantum\_al/circuit.py}'s \texttt{self\_test()}). This confirms the two are genuinely the same mathematical object, not merely analogous constructions -- a necessary but easy-to-get-wrong check, since the naive use of Qiskit's \texttt{SparsePauliOp.from\_operator} silently truncates small coefficients by default and introduced a $\sim\!10^{-5}$ discrepancy until called with explicit zero tolerance.

\subsection{Measurement overhead}

Decomposing a single $K{=}3$ observable bank's variance/covariance terms into Pauli strings, on the padded $32\times32$ (5-qubit) representation, required \textbf{528 distinct Pauli terms for every one of the nine quantities} needed per candidate per iteration ($\mathrm{Var}(\hat{O}_k)$ for $k\in\{\mathrm{struct,elec,thermo}\}$, plus the three symmetrized cross terms) -- out of 1024 possible non-identity Pauli strings on 5 qubits. This is a direct consequence of the observables being dense, block-concentrated matrices rather than sparse or structured operators; naively, a single candidate's full acquisition score would require on the order of $9\times528\approx4750$ measurement settings. This overhead scales combinatorially with the number of observables $K$: the number of covariance pairs is $\binom{K}{2}$ (0 for $K{=}1$, 3 for $K{=}3$, 15 for $K{=}6$), so the sensitivity study's $K{=}6$ configuration (Section~\ref{sec:obssens}) would naively require roughly $21\times528\approx11{,}000$ settings per candidate -- for a configuration that Section~\ref{sec:obssens} already found gives no accuracy benefit over $K{=}3$.

\subsection{Measurement grouping: a real reduction}\label{sec:grouping}

The raw Pauli-term count above does not account for measurement grouping, a standard technique (used throughout the VQE/QAOA resource-estimation literature) that measures multiple commuting Pauli terms simultaneously from a single circuit execution. Applying Qiskit's qubit-wise-commuting (QWC) grouping -- which requires only single-qubit basis-rotation gates before measurement, no additional entangling gates -- to all nine $K{=}3$ quantities reduces the per-quantity term count from 528 to \textbf{122 measurement groups uniformly}, a \textbf{4.3$\times$ reduction} (4752$\to$1098 settings per candidate). Grouping by full (not just qubit-wise) Pauli commutativity reduces this further to \textbf{43 groups per quantity} (387 total, a \textbf{12.3$\times$ reduction}), at the cost of requiring additional Clifford basis-change circuits per group -- a real depth/gate-count trade-off, not accounted for in the group-count alone, so the qubit-wise number is the more honest near-term estimate. This is a genuine, substantial improvement to the feasibility picture: it does not change the shot-count-per-term requirement of Table~\ref{tab:shots} (grouping reduces the number of distinct circuits, not the statistical precision needed from each), but it means the $\sim\!4750$-setting estimate above is a loose upper bound, not the real hardware cost -- roughly 1{,}100 measurement settings per candidate is closer to what a QWC-grouped implementation would actually require.

\subsection{Shot-count convergence of the acquisition ranking}

What matters operationally is not the raw expectation-value error but whether the resulting \emph{ranking} (which candidates get selected) survives finite measurement statistics. Table~\ref{tab:shots} reports, for 30 real band-gap candidates, the mean absolute error against the exact score, Kendall's $\tau$ rank correlation, and the fraction of the exact top-$b$ ($b{=}15$) batch recovered, as a function of shot budget (simulated via Qiskit's precision-scaled estimator, $\mathrm{precision}=1/\sqrt{\mathrm{shots}}$).

\begin{table}[!t]
\centering
\caption{Acquisition-ranking fidelity vs.\ shot budget, 30 real band-gap candidates, noiseless simulation.}
\label{tab:shots}
\begin{tabular}{lccc}
\toprule
Shots & MAE & Kendall $\tau$ & Top-15 overlap \\
\midrule
100 & 0.276 & 0.29 & 67\% \\
1{,}000 & 0.100 & 0.44 & 73\% \\
10{,}000 & 0.024 & 0.77 & 87\% \\
100{,}000 & 0.008 & 0.92 & 100\% \\
\bottomrule
\end{tabular}
\end{table}

The ranking degrades gracefully and recovers fully at high shot counts, confirming the construction is well-posed under realistic finite sampling -- but full batch fidelity requires on the order of $10^5$ shots per Pauli term, which at $\sim\!4750$ terms per candidate is a substantial per-iteration measurement budget on real hardware.

\subsection{Sensitivity to gate noise}

Adding a standard depolarizing noise model (1-qubit gate error $p_1=0.1\%$, 2-qubit/CX error $p_2=1\%$, representative of good current superconducting hardware) at 10,000 shots degrades top-15 overlap from 87\% (noiseless) to 73\%, and Kendall's $\tau$ from 0.77 to 0.63 -- comparable degradation to dropping an order of magnitude in shot budget. The 26-CX, depth-158 state-preparation circuit alone (5 qubits, before any measurement-basis rotations) is enough to accumulate meaningful error at these representative near-term noise rates.

\subsection{Summary}

The formalism admits an exact, verified quantum circuit realization -- this is a genuine (if narrow) confirmation that the mathematics in Section~\ref{sec:method} is not just an analogy. The realization is expensive in exactly the way the formalism's own structure predicts (dense, non-sparse observables; combinatorial growth in $K$; shot budgets further degraded by realistic gate noise), but standard measurement-grouping techniques provide a genuine, substantial ($4$--$12\times$) reduction in that cost, taking the naive $\sim\!4750$-setting-per-candidate estimate down to a more realistic $\sim\!1100$. None of this bears on whether the formalism helps active learning (Sections~\ref{sec:results}--\ref{sec:v3} already answer that, mostly negatively); it characterizes what it would cost, and how that cost can be partially mitigated, to run the classically-simulated experiments above on an actual quantum device -- for a construction that the classical simulation already showed provides no accuracy benefit over simply measuring ensemble disagreement.

\section{Discussion}

\subsection{Why the original formalism does not help, and why the fix that works isn't the quantum part}

The ablation (Table~\ref{tab:ablation}) and the first two improvement attempts (Section~\ref{sec:improvement}) point to a specific diagnosis: \emph{the acquisition score, as originally specified, is computed independently of the downstream predictor}. Definition~\ref{def:total} depends only on $|\psi_i\rangle$, which in turn depends only on candidate $M_i$'s standardized features and the labeled pool's correlation matrix -- never on the residuals, predictions, or disagreement of a model trained on $(X_{\mathrm{train}}, y_{\mathrm{train}})$. Every baseline that beats it does so precisely because it measures \emph{predictive} uncertainty.

Section~\ref{sec:v3}'s residual-coupled variant confirms this diagnosis was correct: once $|\psi_i\rangle$ is built from ensemble disagreement instead of raw features, the method reaches parity with the best baseline on every task. But Table~\ref{tab:v3ablation} shows this improvement has nothing to do with the quantum-inspired apparatus specifically -- the single-observable ablation (essentially a bare quadratic form over per-tree predictions) matches or beats the full covariance-aware, complex-coupled, multi-observable version on four of five tasks, and is dramatically more stable on the fifth. In other words: \emph{what fixed the method was making it more like RF-Uncertainty and QBC, not anything quantum about it.} The covariance coupling, non-commutativity, and complex coefficients that motivate this paper's title and formalism remain, throughout every experiment we ran, either inert or mildly harmful.

\subsection{Relationship to multi-task learning}

Multi-task Gaussian processes and shared-encoder neural networks capture statistical dependence between properties through the \emph{predictive model} architecture \cite{rasmussen2006gaussian}. The formalism studied here instead attempted to capture dependence purely in the \emph{acquisition function}, decoupled from the predictor -- which our results show does not work until it stops being decoupled, at which point the added structure (covariance, complex coupling) provides no further benefit over the disagreement signal alone.

\subsection{Scope and limitations}

(1) \emph{Five real regression tasks, not six}: thermal conductivity was dropped rather than simulated, since it is not broadly available in the Materials Project. (2) \emph{21-dimensional features, not 100}: composition and coarse structural/summary statistics only; no graph-neural-network or structure-aware embeddings. (3) \emph{Tier-1 baselines only}: full multi-output GP and covariance-calibrated deep-ensemble baselines from an earlier version of this manuscript were never actually built and are not compared against here. (4) \emph{Two secondary experiments are unreliable}: the transfer-learning script has an unfixed bug (identical transfer/scratch output) and the multi-property script does not learn effectively in its current form; neither is a basis for any claim in this paper. (5) \emph{Single random-forest predictor}: results may differ with a Gaussian-process or neural-network downstream model, which was not tested here. (6) \emph{Three improvement attempts, not exhaustive search}: we tested three pre-specified, principled modifications (Section~\ref{sec:improvement}, \ref{sec:v3}) and report all of them honestly, stopping once the pattern (ensemble-disagreement coupling helps; quantum-specific machinery does not) was consistent across five tasks; a larger search could in principle find a configuration that wins outright, at increasing risk of fitting noise in a five-task, five-trial evaluation. (7) \emph{Parity, not superiority}: ``statistically indistinguishable from the best baseline'' is not evidence of an advantage -- readers should not read Table~\ref{tab:v3} as the method winning. (8) \emph{Circuit study is a feasibility characterization on a simulator, not hardware}: Section~\ref{sec:hardware}'s shot/noise results use Qiskit Aer's simplified depolarizing model and precision-scaled shot simulation, not an execution on real quantum hardware; real-device results (crosstalk, readout error, T1/T2 decoherence, calibration drift) would very plausibly be worse, not better. The Pauli-decomposition measurement overhead (528 terms) also does not exploit standard variance-reduction techniques (commuting-clique grouping, classical shadows) that would reduce -- though not eliminate -- the measurement count in a real deployment.

\section{Conclusion}

We built and correctly implemented a covariance-aware, quantum-inspired active-learning acquisition function -- state encoding into a feature Hilbert space, non-commuting Hermitian observables, symmetrized cross-observable covariance, and complex coupling coefficients, with a verified classical-limit reduction to standard weighted-variance aggregation. As originally specified, evaluated against nine real baselines on five real Materials Project regression tasks, \textbf{it does not outperform simple baselines}, and none of its constituent mechanisms show an effect distinguishable from noise. Diagnosing the cause -- no dependence on the downstream predictor's residuals -- and fixing it directly (coupling the state encoding to ensemble disagreement) brings the method to statistical parity with the best classical baseline on every task. But a further ablation shows that parity comes entirely from the ensemble-disagreement coupling: the quantum-inspired covariance/multi-observable/complex-coefficient machinery this paper is built around adds no measurable value on top of that signal, and destabilizes it on the noisiest task. The honest conclusion is not ``quantum-inspired uncertainty aggregation helps materials active learning'' but rather: \emph{if you want an acquisition function to benefit from modeling multiple uncertainty sources jointly, first make sure it can see the predictor's own disagreement at all -- the quantum-inspired covariance machinery, at least as formalized here, is not what does the work.}

\balance
\bibliographystyle{IEEEtran}
\bibliography{refs}

\end{document}
